%% 物理第2巻『熱学』：組 熱W9（第22回編集委員会の手引き tools/BRIEF-phys2.md）
%% p2-16.tex  熱-16　読みもの　エネルギーをどう選ぶか（★2〜4。3段にしない）
%% 計画書 tools/phys2-plan.md の B「熱-16」。読みもの（例題は「見本の計算」。じぶんでとこう・章末問題はなし）。第1巻の読みもの p1-16.tex と同じ形。
%% 記号（この章）：効率 η、1秒あたりの熱 Q-dot、電気の出力 P、成績係数 COP_暖（第22回 T-2）、エクセルギー E_x（熱-6 ⑨）、
%%   まわりの温度 T_0〔K〕、ギブズの自由エネルギーの変化 ΔG・エンタルピーの変化 ΔH（熱-10 ⑨）、比熱 c、質量 m。
%%   公式は導き直さず、持ち主の章を指す：η（力-4 ⑧・熱-0 ②）、カルノーの効率（熱-5 ⑤）、COP（熱-5 ②④）、エクセルギー（熱-6 ⑨）、
%%   mc ln(T2/T1)（熱-6 ⑦）、燃焼の原子のバランスと発熱量（熱-1 ⑧）、燃料電池の −ΔG（熱-10 ⑨）、水力 ρgQhη（力-13 ⑪）。
%% 数値はすべて Python で計算した：checks/物理第2巻/熱-16/calc_p2_16.py（出力 calc_p2_16_out.txt）。水の比熱は 4.2 kJ/(kg·K)。
%% ほかの章とそろえた数値：
%%   コンバインド 0.60・石炭火力 0.40・原子力 約0.33（熱-0・熱-5 表 tab-real・熱-8 ⑧）、63.08 %（熱-0・熱-7 ⑧）、
%%   カルノーの効率 45 %（原子力 280 ℃/30 ℃）・65 %（石炭 600 ℃）・84 %（ガス 1600 ℃）（熱-5 表）、
%%   メタンの高発熱量 890.6・低発熱量 802.6 kJ/mol、13A 45 MJ/m^3N（熱-1 ⑧）、ヒートポンプ給湯の COP 3.5（熱-5 章末問題）、
%%   燃料電池 ΔH −285.83・ΔG −237.13 kJ/mol、理論の効率 0.83、1.23 V、0.70 V（熱-10 ⑨・例題）、1 kWh＝3.6 MJ（5-12⑥）。
%% 出典（〔確〕＝検索で確かめた、〔要確認〕＝確かめられなかったので報告に書いた）
%%   ・ライフサイクルの CO2（g-CO2eq/kWh、中央値）：IPCC AR5 第3作業部会 Annex III 表 A.III.2（2014年）
%%     石炭 820、天然ガス（コンバインド）490、太陽光（大規模）48、地熱 38、水力 24、原子力 12、風力（陸上）11。
%%     〔確：二次資料（lowcarbonpower.org の表、Wikipedia の Life-cycle greenhouse gas emissions の表）で照合。原文の PDF は届かなかった。
%%      風力は陸上 11・洋上 12 とする表と、まとめて 11.5 とする表がある〕
%%   ・発電コスト：発電コスト検証ワーキンググループ「基本政策分科会に対する発電コスト検証に関する報告」（2024年12月、資源エネルギー庁）。
%%     2040年に新しくつくる発電所の試算で、太陽光（事業用）7.0〜8.9 円/kWh、原子力 12.5 円/kWh 以上。〔確：報道（2024年12月17日）と
%%     市民団体の解説資料の要約で照合。資源エネルギー庁のサイトには届かなかった〕。ほかの電源の値は紙面に書かない（〔要確認〕）。
%%     太陽光・風力が増えると、出力の変動を補う費用（統合コスト）が加わることも、同じ検証で示された〔確：解説資料の要約〕。
%%   ・設備利用率：同じ 2024 年の検証の前提（太陽光 事業用 18.3 %・住宅 15.8 %、陸上風力 29.6 %、洋上風力 30 %）〔確：検索の要約。原資料は届かなかった〕。
%%     紙面は「太陽光 15〜18 %、風力 約30 %」と丸めて書いた。見本の計算は住宅の屋根の太陽光で 0.15 とした（〔要確認〕：地域・年で変わる）。
%%   ・ITER：2024年7月の新しい計画（ベースライン）で、研究運転の開始 2034年、重水素・三重水素の運転の開始 2039年（ITER 機構の記者会見の要約）〔確〕。
%%     目標は、加熱に入れる 50 MW の10倍（Q＝10）の 500 MW の核融合の出力（ITER 機構の公開資料。広く知られた目標値）。
%%   ・1 toe＝10^7 kcal＝41.868 GJ（IEA・OECD のエネルギー統計の定義。1 kcal＝4.1868 kJ の国際蒸気表カロリー）、1 Btu≈1055 J。
%%   ・太陽電池のモジュールの変換効率 約20 %（市販品のふつうの値）〔要確認：メーカーの公開資料〕。
%%   ・日本の電気の約7割は火力（2020年代のはじめ）〔要確認：資源エネルギー庁「総合エネルギー統計」の確報。サイトに届かなかった〕。
%%   ・日本の発電所の効率は高発熱量で示すことが多い（熱-8 ⑧ の「送電端効率 約42 %（高位発熱量基準）」と同じ）。
\makeatletter
% 横の棒グラフの1本  \qxqbar{y}{ラベル}{長さ(mm)}{色}{値の文字}
\newcommand\qxqbar[5]{%
  \node[font=\sffamily\footnotesize, anchor=east, inner sep=1pt] at (-1.5,#1) {#2};
  \fill[#4] (0,#1-2.2) rectangle (#3,#1+2.2);
  \node[font=\sffamily\footnotesize, anchor=west, inner sep=1pt] at (#3+1.2,#1) {#5};}
\makeatother
\tikzset{qxq lab/.style={font=\sffamily\footnotesize, inner sep=1pt},
  qxq heat/.style={-{Stealth[length=4pt, width=3.6pt]}, line width=1.4pt, bkRed!80!black},
  qxq work/.style={-{Stealth[length=4pt, width=3.6pt]}, line width=1.4pt, bkGreen!55!black}}

\phchapter[読みもの]{16}{エネルギーをどう選ぶか}{電気は、どの方法でつくるとよい？　何をくらべる？}{\bklevelrange{2}{4}}{エネルギーをどう選ぶか}
\bkmichishirube{熱-0,熱-1,熱-5,熱-6,熱-8,熱-9,熱-10,熱-13,熱-15}{電-16,原-15,原-16}

\begin{goal}
\begin{itemize}[label={\color{bkTheme}●}, leftmargin=1.6\zw, itemsep=0pt, topsep=1pt]
\item \bklevel{2}\enspace \qty{1}{kWh}の電気に要る燃料の熱と、燃やして出る二酸化炭素の量を、効率と原子の数のバランスから見積もれる。
\item \bklevel{3}\enspace 発電や給湯の方法を、効率・温度・設備利用率・一生で出る二酸化炭素など、いくつもの物差しでくらべられる。
\item \bklevel{4}\enspace ヒートポンプや燃料電池が得をするわけを、エクセルギー（第二法則）で説明できる。
\end{itemize}
\end{goal}

\begin{hatena}
\noindent\begin{minipage}[c]{0.54\linewidth}
電気は、火力・原子力・水力・太陽光・風力・地熱など、いろいろな方法でつくられています。ニュースでは、「この発電は安い」「二酸化炭素を出さない」「天気で出力が変わる」など、いろいろなことが言われます。\par\vspace{3pt}
電気は、どの方法でつくるとよいのでしょう。そもそも、何と何をくらべればよいのでしょう。
\end{minipage}\hfill
\begin{minipage}[c]{0.42\linewidth}\centering
\begin{tikzpicture}[x=0.9mm, y=0.9mm]
  \fill[bkGray!20] (-2,-2) rectangle (54,0);
  % 火力（建物と煙突）
  \filldraw[fill=bkGray!10, draw=bkInk, line width=0.6pt] (0,0) rectangle (10,8);
  \filldraw[fill=bkGray!10, draw=bkInk, line width=0.6pt] (6,8) rectangle (8,18);
  \fill[bkGray!40] (7,19.5) circle (1.1) (8.4,21.5) circle (1.5);
  % 風車
  \draw[bkInk, line width=0.9pt] (22,0) -- (22,16);
  \foreach \a in {90,210,330} {\draw[bkInk, line width=1.1pt] (22,16) -- ++(\a:7);}
  \fill[bkInk] (22,16) circle (0.8);
  % 太陽と太陽電池
  \fill[bkOrange!80] (44,22) circle (3);
  \foreach \a in {0,45,...,315} {\draw[bkOrange!80, line width=0.6pt] (44,22) ++(\a:4) -- ++(\a:2);}
  \filldraw[fill=bkBlue!40, draw=bkBlue!70!black, line width=0.6pt] (34,1) -- (50,1) -- (46,9) -- (30,9) -- cycle;
  \draw[bkBlue!70!black, line width=0.4pt] (38,1) -- (34,9) (42,1) -- (38,9) (46,1) -- (42,9);
  \draw[bkInk, line width=0.6pt] (40,0) -- (40,1);
  \node[font=\sffamily\bfseries\Large, text=phHatena!80!black] at (13,27) {?};
  \node[font=\sffamily\bfseries\Large, text=phHatena!80!black] at (32,26) {?};
\end{tikzpicture}
\end{minipage}
\end{hatena}

この章は読みものです。じぶんでとこう・章末問題はありません（見本の計算の例題はあります）。これまでの章で学んだ「効率」「温度」「エクセルギー」などの道具を物差しにして、電気や熱のつくり方を、同じ土俵でくらべる方法を見ていきます。先に結論を言うと、「いちばんよい方法」は1つには決まりません。物差しが1つではないからです。でも、物差しをそろえれば、何がわかっていて、何が人の選ぶことなのかを、はっきり分けて話せるようになります。
これまでの章で学んだ道具が、この章でどのように使われるかを、先に見ておきましょう（図\ref{ph:熱-16:fig-michi}）。
\begin{bkfigure}
\begin{tikzpicture}[x=1mm, y=1mm]
  \tikzset{src/.style={bk node, text width=22mm, inner xsep=1mm, align=center, minimum height=17mm, font=\sffamily\footnotesize}}
  \node[src] (a) at (0,24) {効率 $\eta$\\1 kWh の燃料\\（熱-0・熱-8）};
  \node[src] (b) at (27.2,24) {燃える量\\原子のバランス\\（熱-1）};
  \node[src] (c) at (54.4,24) {カルノーの効率\\COP\\（熱-5・熱-9）};
  \node[src] (d) at (81.6,24) {エクセルギー\\$Q(1-T_0/T)$\\（熱-6）};
  \node[src] (e) at (108.8,24) {燃料電池\\$-\Delta G$\\（熱-10）};
  \node[bk node, draw=bkTheme, line width=1pt, text width=96mm, align=left, minimum height=11mm, font=\sffamily\small] (f) at (54.4,-2) {物差しをそろえて、電気や熱のつくり方をくらべる（熱-16）};
  \foreach \n in {a,b,c,d,e} {\draw[bk link] (\n.south) -- (\n.south |- f.north);}
\end{tikzpicture}
\bkfigcaption{これまでの章の道具を、くらべる物差しとして使うのが、この章の方法。}\label{ph:熱-16:fig-michi}
\end{bkfigure}


\begin{yattemiyou}[家の電気の「二酸化炭素の量」を調べよう]
\begin{enumerate}[label={\bkcircled[bkGreen]{\arabic*}}, leftmargin=2.4\zw, itemsep=1pt, topsep=2pt]
\item 家が電気を買っている会社のウェブページで、「電源構成」（どの方法でつくった電気を何割売ったか）を調べる。
\item 同じページで、「CO$_2$ 排出係数」（電気\qty{1}{kWh}あたりに出した二酸化炭素の質量。単位は kg-CO$_2$/kWh）を調べる。何年度の値かも書きとめる。
\item 熱-0 のやってみようで調べた、1か月に使った電気の量（kWh）に、排出係数をかけよう。
\end{enumerate}
\end{yattemiyou}

\noindent たとえば、1か月に\qty{260}{kWh}使い、排出係数が0.45\,kg-CO$_2$/kWh なら、$260\times0.45$ で、約\qty{117}{kg}の二酸化炭素を出したことになります。排出係数は、会社によって、また年によって変わります。その電気を、どの方法でつくったかで変わるからです。この章では、その「どの方法で」の中身を、物理の目で見ていきます。

\begin{omoidasou}
\bkrecall{熱-0}{効率 $\eta$＝電気÷燃料の熱。$\qty{1}{kWh}=\qty{3.6}{MJ}$。電気にならない分は、熱として捨てる。}
\bkrecall{熱-1}{メタンが燃える：$\mathrm{CH_4}+2\mathrm{O_2}\to\mathrm{CO_2}+2\mathrm{H_2O}$。1\,mol で高発熱量\qty{890.6}{kJ}、低発熱量\qty{802.6}{kJ}（⑧）。}
\bkrecall{熱-5}{熱機関の効率の上限は、カルノーの効率 $\eta_{\mathrm C}=1-T_{\mathrm L}/T_{\mathrm H}$（⑤）。ヒートポンプの成績係数は $\mathrm{COP}_{\text{暖}}=Q_{\mathrm H}/W_{\text{入}}$（④）。}
\bkrecall{熱-6}{温度 $T$ の熱 $Q$ のエクセルギー（仕事にできる分）は $Q(1-T_0/T)$。$T_0$ は周りの温度（⑨）。}
\bkrecall{熱-10}{燃料電池が取り出せる電気の仕事は、いちばん多くて $-\Delta G$。水素では理論の効率0.83（⑨）。}
\end{omoidasou}

%% ================================================================
\Needspace{16\baselineskip}
\bksection{doguu}{物理の道具}

\bksubsection{1}{1 kWh の電気に要る燃料の熱}
\leavevmode\label{ph:熱-16:sec1}電気\qty{1}{kWh}は、\qty{3.6}{MJ}のエネルギーです（5-12⑥）。熱機関で電気をつくる発電所では、効率を $\eta$ とすると、\qty{1}{kWh}をつくるのに要る燃料の熱は
\pheq{\text{燃料の熱}=\frac{\qty{3.6}{MJ}}{\eta}}{MJ ＝ MJ ÷ 単位なし}
\label{ph:熱-16:fuel}です（熱-0 ②の効率の式を、燃料の熱について解いたもの）。表\ref{ph:熱-16:tab-fuel}は、熱-0・熱-5・熱-8 で見た効率を入れた値です。
\begin{bkcenter}
\begin{bktable}{colspec={Q[l,m,wd=38mm] Q[c,m,wd=18mm] Q[c,m,wd=30mm] Q[c,m,wd=30mm]}, rowsep=1.2pt, row{1}={bg=bkTheme!10, font=\sffamily\footnotesize}}
発電の方法 & 効率 $\eta$ & 燃料の熱〔MJ〕 & 捨てる熱〔MJ〕 \\
コンバインドサイクル & 約0.60 & 6.0 & 2.4 \\
石炭火力 & 約0.40 & 9.0 & 5.4 \\
原子力 & 約0.33 & 10.9 & 7.3 \\
\end{bktable}
\bktabcaption{電気\qty{1}{kWh}（\qty{3.6}{MJ}）をつくるのに要る燃料の熱と、捨てる熱。捨てる熱は、燃料の熱から\qty{3.6}{MJ}を引いたもの。}\label{ph:熱-16:tab-fuel}
\end{bkcenter}
\begin{phyomi}[燃料の熱 $=\qty{3.6}{MJ}/\eta$]
効率が半分になると、同じ電気に要る燃料は2倍になる。燃料の熱から電気を引いた残りは、熱として周りへ捨てる（熱-0 ④）。効率は「燃料をどれだけ使うか」と「熱をどれだけ捨てるか」を、同時に決めている。
\end{phyomi}
原子力発電所の燃料はウランで、燃やすのではなく、原子核が分かれるとき（核分裂）に出る熱を使います。そのしくみは物理シリーズの『原子』（原-15）で学びます。でも、その熱で水を水蒸気にしてタービンを回す所は、火力と同じランキンサイクル（熱-8）です。だから、蒸気の温度が低い原子力の効率は、火力より低くなります（熱-5 ⑦）。

\bksubsection{2}{燃やすと出る二酸化炭素：原子の数のバランス}
\leavevmode\label{ph:熱-16:sec2}燃料を燃やすと、燃料の中の炭素 C は、すべて二酸化炭素 CO$_2$ になります。燃える前と後で、原子の数は変わりません（熱-1 ⑧の原子のバランス）。
\[\mathrm{CH_4}+2\mathrm{O_2}\to\mathrm{CO_2}+2\mathrm{H_2O}\qquad\qquad \mathrm{C}+\mathrm{O_2}\to\mathrm{CO_2}\]
メタン（天然ガスのおもな成分）\qty{1}{mol}が燃えると、CO$_2$ が\qty{1}{mol}（\qty{44}{g}）でき、低発熱量で\qty{802.6}{kJ}の熱が出ます。炭素（石炭のおもな成分）\qty{1}{mol}が燃えても、CO$_2$ は\qty{1}{mol}（\qty{44}{g}）ですが、出る熱は\qty{393.5}{kJ}だけです。熱\qty{1}{MJ}あたりにすると、
\[\text{メタン}：\ \frac{44\times1000}{802.6}\approx\qty{55}{g/MJ},\qquad \text{炭素}：\ \frac{44\times1000}{393.5}\approx\qty{112}{g/MJ}\]
で、炭素はメタンの約2倍の CO$_2$ を出します。メタンは、炭素のほかに水素も燃えて熱を出し、水素が燃えてできるのは水だからです。

\Needspace{6\baselineskip}
\phboxhead{1 kWh あたりにすると}
1 kWh の燃料の熱（①）に、1 MJ あたりの CO$_2$ をかければ、1 kWh あたりの CO$_2$ になります。
\pheq{\text{1 kWh の CO$_2$}=\frac{\qty{3.6}{MJ}}{\eta}\times\text{（1 MJ あたりの CO$_2$）}}{g ＝ MJ × g/MJ}
\label{ph:熱-16:co2}コンバインドサイクル（$\eta=0.60$）なら $6.0\times55\approx\qty{330}{g}$、石炭火力（$\eta=0.40$、石炭を炭素とみる）なら $9.0\times112\approx\qty{1000}{g}$ で、約3倍の差です。詳しい計算は\bkreidairef{1}でします。
\begin{phbaisuru}[0.40から0.60に]{効率}
同じ燃料なら、1 kWh あたりの CO$_2$ は効率に反比例するので、$0.40\div0.60=\tfrac23$ 倍になります。燃料をメタンに変えると、さらに約半分です。CO$_2$ を減らす方法には、「燃料を変える」と「効率を上げる」の2つがあり、この式はその両方を表しています。CO$_2$ を出さない方法（③）にすれば、燃やす所から出る CO$_2$ は0になります。
\end{phbaisuru}

\Needspace{8\baselineskip}
\phboxhead{表計算で、くらべの表をつくる}
①と②の計算を表計算にすると、効率や燃料を変えたときの CO$_2$ を、すぐにくらべられます。
\begin{bkcenter}
\bksheet{5}{21}{%
1 & 方法 & 効率 & 燃料の熱（MJ） & CO$_2$（g/MJ） & CO$_2$（g/kWh） \\
2 & ガス・コンバインド & 0.60 & 6.0 & 54.8 & 329 \\
3 & ガス・汽力 & 0.40 & 9.0 & 54.8 & 493 \\
4 & 石炭（炭素） & 0.40 & 9.0 & 111.8 & 1006 \\
5 & 石炭（炭素） & 0.45 & 8.0 & 111.8 & 895 \\}
\bktabcaption{C2 は \texttt{=3.6/B2}、E2 は \texttt{=C2*D2}。D2 はメタンの \texttt{=44000/802.6}、D4 は炭素の \texttt{=44000/393.5}。「汽力」は、ボイラーで燃料を燃やして蒸気だけでタービンを回す火力。}\label{ph:熱-16:tab-sheet}
\end{bkcenter}
\noindent 同じ天然ガスでも、コンバインドサイクルにすると、汽力より CO$_2$ が3分の1減ります（2行め・3行め）。石炭の効率を0.40から0.45に上げても、ガスのコンバインドには遠くおよびません（4行め・5行め）。B列の効率を書きかえて、どこまで上げればガスに追いつくか、ためしてみましょう。

\bksubsection{3}{発電の方法をくらべる：何を電気にしているか}
\leavevmode\label{ph:熱-16:sec3}発電の方法によって、「電気にしているもの」がちがいます。だから、効率の分母（入れたもの）もちがいます。
\begin{bkcenter}
\small
\begin{tblr}{colspec={Q[l,m,wd=22mm] Q[l,m,wd=30mm] X[l,m] Q[c,m,wd=14mm]}, width=\linewidth, rowsep=1.1pt, colsep=3pt,
  hlines={0.4pt,bkLine}, row{1}={bg=bkTheme!12, font=\sffamily\bfseries\footnotesize}}
方法 & 電気にしているもの & 上限を決めるもの & 章 \\
火力・原子力・地熱 & 熱（蒸気・ガス） & 熱を受けとる温度と捨てる温度（カルノーの効率） & 熱-5・8 \\
水力 & 水の位置エネルギー & 熱を通らないので、カルノーの上限はない。水車と発電機の損失だけ & 力-13 \\
風力 & 風の運動エネルギー & 風を止めきれない（もっと知りたい） & 力-13 \\
太陽光 & 太陽の光 & 太陽電池の材料と、光の波長（色） & 熱-13 \\
燃料電池 & 燃料の化学エネルギー & $-\Delta G$（熱を通らないので、カルノーの上限はない） & 熱-10 \\
\end{tblr}
\bktabcaption{発電の方法と、電気にしているもの。}\label{ph:熱-16:tab-method}
\end{bkcenter}

\phboxhead{熱で電気をつくる方法：温度が高いほど有利}
火力・原子力・地熱は、熱でタービンを回します。熱の中で仕事にできる分（エクセルギー）の割合は、周りの温度を $T_0$ とすると $1-T_0/T$ でした（熱-6 ⑨）。熱を受けとる温度 $T$ が高いほど、この割合が大きくなります（図\ref{ph:熱-16:fig-exergy}）。
\begin{bkfigure}
\begin{tikzpicture}
\begin{axis}[bkplot, width=110mm, height=42mm, xmin=0, xmax=1750, ymin=0, ymax=1, xtick={300,600,900,1200,1500}, ytick={0.2,0.4,0.6,0.8,1}, xlabel={熱を受けとる温度（℃）}, ylabel={$1-T_0/T$}]
  \addplot[bkcurve, domain=30:1720, samples=80] {1-303.15/(x+273.15)};
  \fill[bkRed] (axis cs:280,0.452) circle (2pt) (axis cs:600,0.653) circle (2pt) (axis cs:1600,0.838) circle (2pt);
  \node[font=\sffamily\footnotesize, anchor=north west] at (axis cs:300,0.40) {原子力の蒸気 約280\,℃：0.45};
  \node[font=\sffamily\footnotesize, anchor=north west] at (axis cs:620,0.60) {石炭火力の蒸気 約600\,℃：0.65};
  \node[font=\sffamily\footnotesize, anchor=south east] at (axis cs:1640,0.88) {ガスタービン 約1600\,℃：0.84};
\end{axis}
\end{tikzpicture}
\bkfigcaption{熱のうち仕事にできる割合 $1-T_0/T$（周りは $T_0=\qty{303}{K}$、約30\,℃）。熱-5 の表のカルノーの効率と同じ数。}\label{ph:熱-16:fig-exergy}
\end{bkfigure}
コンバインドサイクルが効率のよい理由の1つは、ガスタービンで1600\,℃級の高い温度の熱を受けとり、蒸気タービンで約30\,℃の海水へ熱を捨てるからです（熱-7 ⑧）。反対に、地熱や海洋温度差発電のように温度の低い熱では、どんなにうまく作っても、効率は低くなります（熱-5 ⑦）。

\phboxhead{熱を通らない方法：効率の意味がちがう}
水力・風力・太陽光は、燃料を燃やしません。入れるもの（水・風・光）は、使ってもなくならず、お金もかかりません。だから、効率が低くても燃料をむだにしているわけではありません。効率が低いと、同じ電気をつくるのに、大きなダム・大きな羽根・広い面積が要る、ということです。たとえば、太陽電池の効率は約2割で、地面に届く日光を\qty{1}{m^2}あたり約\qty{1000}{W}（よく晴れた昼）とすると、\qty{1}{kW}をつくるには約\qty{5}{m^2}の太陽電池が要ります（\bkreidairef{3}）。
\begin{phmijika}[ダムの水と、川の流れ]
水力発電所の出力は $\rho gQh\eta$（力-13 ⑪）で、効率 $\eta$ は0.8〜0.9と、熱機関よりずっと高くなります。位置エネルギーを、熱を通らずに、水車で直接回す仕事に変えているからです。そのかわり、落差 $h$ と流量 $Q$ は、山と川の形で決まり、どこにでもつくれるわけではありません。
\end{phmijika}

\phboxhead{エネルギーは減らないのに、なぜ「使う」と言う？}
エネルギーは保存します（熱-2）。発電所の燃料の熱も、家で使った電気も、最後は熱として周りの空気や海を少し温め、宇宙へ放射で出ていきます（熱-13）。エネルギーの量は、どこでも減っていません。それでも「エネルギーを使う」と言うのは、仕事にできる分、つまりエクセルギーが減るからです（熱-6 ⑨）。高い温度の炎、高い所の水、強い風は、エクセルギーの大きい、そろった状態です。使ったあとの、周りと同じ温度のぬるい熱は、分子の動きがばらばらに散らばった状態で、並び方の数がとても多い状態です（熱-15）。わたしたちが選んでいるのは、どのエクセルギーを、どれだけむだなく、どの順に使うか、なのです。

\begin{motto}[3]
\phboxhead{風から取り出せる力の上限}
風の運動エネルギーが1秒に羽根の円を通る量は、空気の密度 $\rho$、羽根の円の面積 $A$、風の速さ $v$ を使って
\[\frac{1}{2}\rho Av^3\]
です（1秒に通る空気の質量 $\rho Av$ に、1 kg あたりの運動エネルギー $v^2/2$ をかけた。力-13）。羽根の直径\qty{100}{m}、風の速さ\qty{10}{m/s}なら約\qty{4.7}{MW}、風の速さが半分の\qty{5}{m/s}なら $\tfrac18$ の約\qty{0.59}{MW}です。風を完全に止めると、空気が羽根を通りぬけられなくなるので、全部は取り出せません。運動量とエネルギーのバランス式（力-13）から、取り出せるのは多くても $\tfrac{16}{27}$（約59\,\%）とわかっています（ベッツの限界。導き方は専門の本で扱います）。
\end{motto}

\bksubsection{4}{一生で数える二酸化炭素：ライフサイクル}
\leavevmode\label{ph:熱-16:sec4}②で数えたのは、燃料を燃やす所で出る CO$_2$ だけでした。でも、発電所をつくるための鉄やコンクリート、燃料を掘る・運ぶ、使い終わった設備をこわす、といった所でも、CO$_2$ は出ます。発電所の一生（ライフサイクル）全体で出る量を、一生でつくる電気の量でわったものを、\textbf{ライフサイクルの CO$_2$} といいます。\bkidx{らいふさいくる}{ライフサイクル}このように、物の一生を通して環境への影響を数える方法を、ライフサイクルアセスメントといいます。

図\ref{ph:熱-16:fig-lca}は、世界のたくさんの研究の値を集めて、国連の気候変動に関する政府間パネル（IPCC）が2014年の報告書にまとめた中央値です。CO$_2$ のほかに、メタンなどの温室効果ガスも、CO$_2$ に直して数えてあります。
\begin{bkfigure}
\begin{tikzpicture}[x=1mm, y=1mm]
  \qxqbar{42}{石炭}{82}{bkGray!70}{820}
  \qxqbar{35}{天然ガス（コンバインド）}{49}{bkOrange!70}{490}
  \qxqbar{28}{太陽光（大規模）}{4.8}{bkYellow!80!black}{48}
  \qxqbar{21}{地熱}{3.8}{bkRed!60}{38}
  \qxqbar{14}{水力}{2.4}{bkBlue!60}{24}
  \qxqbar{7}{原子力}{1.2}{bkPurple!60}{12}
  \qxqbar{0}{風力（陸上）}{1.1}{bkGreen!60}{11}
  \draw[bkInk, line width=0.5pt] (0,-4) -- (0,46);
  \draw[bkInk, line width=0.5pt] (0,-4) -- (90,-4);
  \foreach \x/\t in {0/0, 20/200, 40/400, 60/600, 80/800} {\draw[bkInk, line width=0.5pt] (\x,-4) -- (\x,-5.2); \node[qxq lab, anchor=north] at (\x,-5.6) {\t};}
  \node[qxq lab, anchor=north] at (45,-10.5) {1 kWh あたりの CO$_2$（g。ほかの温室効果ガスも CO$_2$ に直してふくむ）};
\end{tikzpicture}
\bkfigcaption{発電の方法ごとの、一生で出る CO$_2$（IPCC 第5次評価報告書・第3作業部会（2014年）の中央値）。同じ方法でも、研究によって何割も上下する。}\label{ph:熱-16:fig-lca}
\end{bkfigure}
\noindent 燃料を燃やす火力は、1 kWh あたり数百 g です。燃やさない方法は数十 g 以下で、火力の数十分の1です。②の見積もり（コンバインド 約\qty{330}{g}、石炭 約\qty{1000}{g}）は、燃やす所だけを数え、燃料を純粋なメタン・炭素とみたものでした。図の値は、世界のいろいろな効率の発電所、燃料の成分、掘る・運ぶ所でもれるメタンなどもふくめた値です。どちらも、火力は数百 g、石炭はガスの約2〜3倍、という桁は同じです。

\bksubsection{5}{出したいときに出せるか：設備利用率}
\leavevmode\label{ph:熱-16:sec5}太陽光は夜には電気をつくれず、風力は風が弱いと止まります。発電所が1年間に実際につくった電気の量を、ずっと定格の出力で動いたときの量でわった割合を、\textbf{設備利用率}といいます。\bkidx{せつびりようりつ}{設備利用率}
\pheq{\text{設備利用率}=\frac{\text{1年につくった電気〔kWh〕}}{\text{定格の出力〔kW〕}\times\qty{8760}{h}}}{単位なし}
1年は $365\times24=\qty{8760}{h}$ です。国の2024年の試算では、太陽光は15〜18\,\%、風力は約30\,\%と見こんでいます。燃料を入れれば動く火力や原子力は、点検の間などを除けば、ずっと動かすことができます。

\phboxhead{電気は、つくる量と使う量を合わせ続ける}
電気は、つくったその瞬間に使われます。つくる量と使う量がずれると、電気の周波数（東日本で 50\,Hz、西日本で 60\,Hz）が変わり、ずれが大きいと停電します。そこで、天気で変わる太陽光や風力を多く使うときには、出力をすばやく変えられる火力、ためておく設備、遠くの地域と電気をやりとりする送電線などを組み合わせます。そのしくみは、物理シリーズの『電磁気学』（電-16）で学びます。

\phboxhead{電気をためる：揚水発電}
電気の余る時間に、下の池の水をポンプで上の池へくみ上げ、電気の足りない時間に、その水を落として発電するのが揚水発電です。水\qty{1}{m^3}（\qty{1000}{kg}）を\qty{500}{m}上げると、位置エネルギーは $1000\times9.8\times500=\qty{4.9e6}{J}$、約\qty{1.4}{kWh}です（力-4）。100万 kW を6時間出し続けるには、くみ上げとくみ出しの損失を無視しても、約440万 m$^3$ の水が要ります。熱-0 の100万 kW の発電所が1日に使う海の水（約450万 m$^3$）と同じくらいの量です。

\bksubsection{6}{熱がほしいときは：ヒートポンプと給湯}
\leavevmode\label{ph:熱-16:sec6}家で使うエネルギーのうち、お湯をわかす熱と、部屋を暖める熱は、大きな割合をしめます。熱をつくる方法も、物差しをそろえてくらべてみましょう。\qty{300}{L}の水を15\,℃から45\,℃にする熱は、$300\times4.2\times30=\qty{37800}{kJ}$、つまり\qty{10.5}{kWh}です。これを、3つの方法でつくります（\bkreidairef{2}）。
\begin{enumerate}[label={\bkcircled[bkBlue]{\arabic*}}, leftmargin=2.4\zw, itemsep=1pt, topsep=2pt]
\item \textbf{ガスの給湯器}（効率0.95）：ガスの熱が\qty{11}{kWh}要る。
\item \textbf{電気のヒーター}：電気が\qty{10.5}{kWh}要る（電気の仕事を、全部そのまま熱にする）。
\item \textbf{ヒートポンプの給湯器}（$\mathrm{COP}_{\text{暖}}=3.5$、熱-5 ④・熱-9）：電気が\qty{3.0}{kWh}要る。残りの\qty{7.5}{kWh}は、外の空気からくみ上げた熱。
\end{enumerate}
電気は、発電所で燃料からつくっています。効率0.60の発電所でつくるとすると、電気のヒーターは燃料の熱\qty{17.5}{kWh}、ヒートポンプは\qty{5.0}{kWh}を使うことになります（図\ref{ph:熱-16:fig-yu}）。
\begin{bkfigure}
\begin{tikzpicture}[x=1mm, y=1mm]
  \qxqbar{24}{ガスの給湯器}{55.3}{bkOrange!70}{11.1}
  \qxqbar{16}{電気のヒーター}{87.5}{bkRed!55}{17.5}
  \qxqbar{8}{ヒートポンプ}{25}{bkBlue!55}{5.0}
  \qxqbar{0}{理想（エクセルギーの分）}{4.3}{bkGreen!60}{0.85}
  \draw[bkInk, line width=0.5pt] (0,-4) -- (0,28);
  \draw[bkInk, line width=0.5pt] (0,-4) -- (95,-4);
  \foreach \x/\t in {0/0, 25/5, 50/10, 75/15} {\draw[bkInk, line width=0.5pt] (\x,-4) -- (\x,-5.2); \node[qxq lab, anchor=north] at (\x,-5.6) {\t};}
  \node[qxq lab, anchor=north] at (45,-10.5) {使う燃料の熱（kWh。電気は効率0.60の発電所でつくるとした）};
\end{tikzpicture}
\bkfigcaption{\qty{300}{L}の水を15\,℃から45\,℃にする（\qty{10.5}{kWh}の熱）のに使う燃料の熱。いちばん下は、第二法則が許すいちばん少ない量。}\label{ph:熱-16:fig-yu}
\end{bkfigure}
\noindent ヒートポンプは、発電所で熱を捨てても、なおガスの給湯器の半分以下の燃料ですみます。\bkidx{ひーとぽんぷ}{ヒートポンプ}これは、熱-0 の熱の流れの帯では見えない「熱の質」のちがいによるものです。

\phboxhead{45\,℃のお湯は「質の低い」熱}
45\,℃のお湯の熱は、周り（15\,℃）との温度差が小さいので、エクセルギー（仕事にできる分）がとても小さい熱です。お湯の温度は15\,℃から45\,℃まで変わっていくので、平均の温度約30\,℃（\qty{303}{K}）の熱とみなして、熱-6 ⑨の式に入れると、
\[E_{\mathrm x}\approx Q\Bigl(1-\frac{T_0}{T}\Bigr)=10.5\times\Bigl(1-\frac{288}{303}\Bigr)\approx\qty{0.52}{kWh}\]
です。お湯をわかすのに本当に要る仕事は、熱\qty{10.5}{kWh}のうちの約\qty{0.5}{kWh}だけなのです。ガスの炎（1000\,℃をこえる）や電気（エクセルギーそのもの）で45\,℃のお湯をつくると、残りのエクセルギーは、温度差のある所で熱が流れるときに失われます（熱-6 ⑨の $T_0S_{\mathrm{gen}}$）。ヒートポンプは、外の空気の熱を使って、その失われる分を減らしているのです。

\begin{motto}[4]
\phboxhead{お湯に要るいちばん少ない仕事}
水 $m$ を $T_0$ から $T$ まで温めるとき、可逆なヒートポンプが入れる仕事が、いちばん少なくなります。そのとき、外の空気から熱を受けとって、全体で生まれるエントロピーは0です。水のエントロピーは $mc\ln(T/T_0)$ 増える（熱-6 ⑦）ので、外の空気から受けとる熱は $T_0\,mc\ln(T/T_0)$ で、仕事はその残りです。
\[W_{\min}=mc\,(T-T_0)-T_0\,mc\ln\frac{T}{T_0}\]
\qty{300}{kg}、$T_0=\qty{288.15}{K}$、$T=\qty{318.15}{K}$ では $W_{\min}\approx\qty{1.84e3}{kJ}$（\qty{0.51}{kWh}）で、本文の見積もり（\qty{0.52}{kWh}）とほぼ合います。理想の成績係数は $10.5\div0.51\approx20$ です。本物のヒートポンプ（3〜4）がこれより低いのは、熱交換器で温度差をつけて熱を移すこと（熱-13）や、圧縮機の損失のためです。
\end{motto}

\bksubsection{7}{燃料電池と水素}
\leavevmode\label{ph:熱-16:sec7}燃料電池は、水素と酸素から水をつくる反応のエネルギーを、燃やさずに、直接電気にします（熱-10 ⑨）。\bkidx{ねんりょうでんち}{燃料電池}熱をいったん高い温度にしてから仕事に変える熱機関ではないので、カルノーの効率の上限は当てはまりません。理論の効率は、反応の熱 $-\Delta H$ に対する電気の仕事 $-\Delta G$ の割合で、25\,℃で液体の水ができるとき $237.13\div285.83\approx0.83$ です（\bkreidairef{4}）。

ただし、水素は、天然ガスや石炭のように地面の下にたまっているわけではありません。水を電気で分解（電気分解）するか、天然ガスなどから取り出してつくります。水を分解するには、少なくとも、燃料電池で取り出せる電気の仕事と同じだけ（\qty{1}{mol}あたり\qty{237}{kJ}）の電気が要ります。だから、水素は、電気や熱を「生み出す」ものではなく、エネルギーを「運ぶ・ためる」ものと考えます。電気→水素→電気ともどすと、それぞれの段の損失の分だけ、電気は減ります。天然ガスから水素をつくると、そのとき CO$_2$ が出ます。
\begin{phmijika}[燃料電池の車と、家の燃料電池]
燃料電池の車は、タンクの水素で電気をつくり、モーターで走ります。走るときに出るのは水だけです。家庭用の燃料電池は、都市ガスから水素を取り出して電気をつくり、そのとき出る熱でお湯もわかします。電気と熱の両方を使うので、燃料の熱を、むだなく使えます（工学ではこう使う）。
\end{phmijika}

\Needspace{10\baselineskip}
\bksubsection{8}{統計の数を読む：何を、どの基準で数えたか}
\leavevmode\label{ph:熱-16:sec8}エネルギーの統計やニュースの数を読むときは、次の3つに気をつけます。

\phboxhead{一次エネルギーと最終エネルギー}
石油・石炭・天然ガス・ウラン・水の落ちる力・日光のように、自然から取り出したままの形のエネルギーを\textbf{一次エネルギー}といいます。\bkidx{いちじえねるぎー}{一次エネルギー}家・工場・車で使う電気・ガソリン・都市ガスなどの形を、\textbf{最終エネルギー}といいます。\bkidx{さいしゅうえねるぎー}{最終エネルギー}一次エネルギーから最終エネルギーに変える所（おもに発電所）で、熱-0 で見た「捨てる熱」が出るので、最終エネルギーは一次エネルギーよりかなり小さくなります。

\phboxhead{高発熱量か、低発熱量か}
燃料の熱を、できた水の蒸発熱までふくめて数えるのが高発熱量、ふくめないのが低発熱量でした（熱-1 ⑧）。メタンでは、低発熱量は高発熱量の $802.6\div890.6\approx0.90$ 倍です。同じ発電所でも、燃料の熱を高発熱量で数えると分母が大きくなるので、効率は小さく出ます。たとえば、63.08\,\%（熱-0）が低発熱量で数えた効率なら、高発熱量で数えると
\[0.6308\times\frac{802.6}{890.6}\approx0.57\]
です。日本の発電所の効率は高発熱量で示すことが多く（熱-8 ⑧）、ガスタービンなどの機械の性能は低発熱量で示すことがあります。効率の数をくらべるときは、どちらの基準かを必ず確かめます。

\begin{senmonsho}[発電設備の仕様の書き方]
\begin{bkhonbun}
本設備は一軸型のコンバインドサイクル発電設備であり，定格出力 500 MW，発電端熱効率 63\,\%（低位発熱量基準）である．所内率は 2\,\% とする．
\end{bkhonbun}
\bkyaku{「発電端」は発電機の出口で測った電気、「送電端」は、発電所の中で使う電気（所内電力）を引いて、送電線へ送り出す電気。「所内率」は、発電した電気のうち、発電所の中で使う割合。「一軸型」は、ガスタービン・蒸気タービン・発電機を1本の軸でつないだ形。「低位発熱量」は、この本の低発熱量（熱-1 ⑧）と同じ。効率の数を読むときは、「発電端か送電端か」「高位か低位か」の2つを確かめる（\bkreidairef{5}）。}
\end{senmonsho}

\phboxhead{何年の、だれの数か}
電源構成・費用・排出係数は、年によって変わります。数を使うときは、「何年の、どの統計（だれが出した数）か」を必ず書きます。この本でも、図や本文の数には、年と出典を書いてあります。数が少し変わっても結論が変わらないかどうかを、確かめておくことも大切です。たとえば、図\ref{ph:熱-16:fig-lca}の値が研究によって何割か変わっても、「火力は数百 g、燃やさない方法は数十 g 以下」という桁のちがいは変わりません。

\bksubsection{9}{はてな？の答え：物差しをそろえて、くらべる}
\leavevmode\label{ph:熱-16:sec9}\bkhatenaref に答えましょう。電気のつくり方をくらべる物差しは、少なくとも表\ref{ph:熱-16:tab-monosashi}の6つがあります。
\begin{bkcenter}
\small
\begin{tblr}{colspec={Q[l,m,wd=26mm] X[l,m] Q[l,m,wd=36mm]}, width=\linewidth, rowsep=1.1pt, colsep=3pt,
  hlines={0.4pt,bkLine}, row{1}={bg=bkTheme!12, font=\sffamily\bfseries\footnotesize}}
物差し & 何で測るか（この本の道具） & 気をつけること \\
燃料の量 & 効率、1 kWh の燃料の熱 $\qty{3.6}{MJ}/\eta$（①） & 分母は何か。発熱量の基準 \\
二酸化炭素 & 1 kWh あたりの CO$_2$（②・④） & 燃やす所だけか、一生か \\
費用 & 1 kWh あたりの円（発電コスト） & 何年の試算か。前提は何か \\
出したいときに & 設備利用率、ためる・合わせる（⑤） & 夜・天気・季節 \\
資源 & 燃料はどこから、いつまで手に入るか & 輸入にたよるか、ためておけるか \\
安全・環境 & 事故への備え、使い終わった燃料、温排水、景観 & 物理だけでは決まらない \\
\end{tblr}
\bktabcaption{電気のつくり方をくらべる物差し。}\label{ph:熱-16:tab-monosashi}
\end{bkcenter}
\noindent 費用も、年と前提で大きく変わります。国の2024年の試算では、2040年に新しくつくる発電所の費用は、太陽光（事業用）が1 kWh あたり7.0〜8.9円、原子力が12.5円以上と見積もられました。ただし、太陽光や風力を多く使うと、天気で変わる出力を補う費用（\textbf{統合コスト}）が加わることも、同じ試算で示されています。\bkidx{とうごうこすと}{統合コスト}

どの物差しでも1番になる方法は、ありません。燃やす火力は、出したいときに出せますが、CO$_2$ を出し、燃料の多くを輸入しています。太陽光と風力は、燃料が要らず CO$_2$ も少ないのですが、天気で出力が変わります。原子力は、CO$_2$ が少なく、ずっと動かせますが、事故への備えと、使い終わった燃料の扱いが大きな課題です（原-16）。水力・地熱は、つくれる場所が限られます。だから、多くの国は、いくつもの方法を組み合わせて電気をつくっています。どれをどれだけ組み合わせるかは、物差しのどれを重く見るかで変わる、社会の選択です。この本で学んだ物理は、その選択の材料になる数を、正しく出すための道具なのです。

\begin{phmada}[核融合で、電気をつくれるか？]
太陽は、水素の原子核どうしがくっつく\textbf{核融合}のエネルギーで光っています（原-11）。地上で核融合を起こして電気をつくる研究が、世界で続いています。日本もふくむ国々が南フランスで建設している実験炉 ITER は、燃料を温めるために外から入れるエネルギー（1秒あたり\qty{50}{MW}）の10倍の、核融合の出力（\qty{500}{MW}）を出すことをめざしています。2024年に見直された計画では、研究のための運転を2034年に始め、燃料に重水素と三重水素を使う運転は2039年に始める予定です。ITER は電気をつくる発電所ではありません。核融合の熱を、長い期間、安く、安全に電気に変えられるかどうかは、そのあとの研究で確かめることで、まだわかっていません。
\end{phmada}
\begin{phmada}[たくさんの電気を、安く、長くためられるか？]
揚水発電（⑤）や電池は、数時間から1日くらいの電気をためるのに使われています。でも、季節をこえるような長い期間、国全体を動かせるほどの電気を、安くためる方法は、まだ見つかっていません。水素にして貯める（⑦）、熱にして貯める、などの方法が研究されていますが、どの方法が、どれだけ安くなるかは、まだわかっていません。
\end{phmada}
\begin{phsaki}{エネルギー工学と電力系統}
発電所の熱の流れとエクセルギーを、設備の1つ1つまで数えて設計するのは、\textbf{エネルギー工学}です。発電所・送電線・ためる設備をつないで、つくる量と使う量をいつも合わせる\textbf{電力系統}のしくみは、物理シリーズの『電磁気学』（電-16）で学びます。核分裂の原子炉のしくみと、使い終わった燃料の扱いは、『原子』（原-15・原-16）で学びます。
\end{phsaki}

\begin{kotoba}\small
\bkword{一次エネルギー}{石油・石炭・天然ガス・ウラン・水力・日光など、自然から取り出したままの形のエネルギー。}
\bkword{最終エネルギー}{電気・ガソリン・都市ガスなど、家・工場・車で使う形のエネルギー。}
\bkword{設備利用率}{1年につくった電気を、ずっと定格の出力で動いたときの量でわった割合。}
\bkword{ライフサイクル}{発電所をつくってから、こわすまでの一生。その全体で出る CO$_2$ を数える。}
\bkword{統合コスト}{天気で変わる出力を補うための、ほかの設備や送電線の費用。}
\end{kotoba}

%% ================================================================
\Needspace{8\baselineskip}\bksection{reidai}{見本の計算}

\begin{phreidai}[1 kWh の電気で出る二酸化炭素]\label{ph:熱-16:ex1}
電気\qty{1}{kWh}を、（1）天然ガスのコンバインドサイクル（効率0.60、低発熱量で数えた効率）と、（2）石炭火力（効率0.40）でつくります。燃やす所で出る CO$_2$ の質量を求めましょう。天然ガスはメタン（低発熱量\qty{802.6}{kJ/mol}）、石炭は炭素（燃える熱\qty{393.5}{kJ/mol}）とします。
\phbamen
同じ\qty{1}{kWh}の電気でも、燃料と効率で CO$_2$ が変わります。
\phmushi
燃料はすべて燃えて、炭素は全部 CO$_2$ になるとします。燃料を掘る・運ぶ所や、発電所をつくる所で出る CO$_2$ は数えません。石炭の中の水素や灰は無視します。
\phzu
効率 $\eta$、燃料の熱 $Q=\qty{3.6}{MJ}/\eta$、CO$_2$ のモル質量 \qty{44}{g/mol}。
\phriyu
燃料の熱は①の式で出ます。CO$_2$ の物質量は、燃料の物質量と同じ（原子のバランス、熱-1 ⑧）で、燃料の物質量は「熱÷1 mol あたりの熱」です。
\phtateru
\[\text{CO$_2$ の質量}=\frac{Q}{\text{1 mol あたりの熱}}\times\qty{44}{g/mol}\]
\phtoku
（1）\[Q=\frac{3.6}{0.60}=\qty{6.0}{MJ},\qquad \frac{6000}{802.6}\times44\approx7.48\times44\approx\qty{329}{g}\]
（2）\[Q=\frac{3.6}{0.40}=\qty{9.0}{MJ},\qquad \frac{9000}{393.5}\times44\approx22.9\times44\approx\qty{1.01e3}{g}\]
\phtashikame
単位：kJ÷（kJ/mol）×（g/mol）＝g。\\
大きさ：（2）は（1）の約3倍です。そのうち2倍は燃料のちがい（②の\qty{112}{g/MJ}と\qty{55}{g/MJ}）、1.5倍は効率のちがい（$0.60\div0.40$）です。図\ref{ph:熱-16:fig-lca}の一生の値（490・820）とも、桁が合っています。\\
極端な場合：効率が1でも、（1）は $3.6\times55\approx\qty{200}{g}$ より少なくはなりません。燃料を変えずに CO$_2$ を0にはできません。
\kotae （1）約\qty{330}{g}（2）約\qty{1000}{g}
\end{phreidai}

\begin{phreidai}[300 L のお湯を3つの方法でわかす]\label{ph:熱-16:ex2}
\qty{300}{L}の水を15\,℃から45\,℃にします。（1）お湯に入れる熱、（2）ヒートポンプの給湯器（$\mathrm{COP}_{\text{暖}}=3.5$）が使う電気、（3）ガスの給湯器（効率0.95）が使うガスの熱と、都市ガス（13A、\qty{45}{MJ/m^3}）の体積を求めましょう。
\phbamen
1日に使うお湯を、夜のうちにまとめてわかします。
\phmushi
水の比熱は一定、水\qty{1}{L}は\qty{1}{kg}とします。タンクや配管から逃げる熱は、成績係数・効率の中にふくめます。
\phzu
$m=\qty{300}{kg}$、$c=\qty{4.2}{kJ/(kg.K)}$、$\Delta T=45-15=\qty{30}{K}$。お湯に入れる熱 $Q$、ヒートポンプの電気 $W_{\text{入}}$。
\phriyu
お湯の熱は $Q=mc\,\Delta T$（熱-1 ③）。ヒートポンプは $\mathrm{COP}_{\text{暖}}=Q/W_{\text{入}}$（熱-5 ②）、ガスは効率の式（力-4 ⑧）。
\phtateru
\[Q=mc\,\Delta T,\qquad W_{\text{入}}=\frac{Q}{\mathrm{COP}_{\text{暖}}},\qquad \text{ガスの熱}=\frac{Q}{0.95}\]
\phtoku
（1）\[Q=300\times4.2\times30=\qty{37800}{kJ}=\frac{37800}{3600}\,\text{kWh}=\qty{10.5}{kWh}\]
（2）\[W_{\text{入}}=\frac{10.5}{3.5}=\qty{3.0}{kWh}\]
（3）\[\frac{37800}{0.95}\approx\qty{39800}{kJ}\approx\qty{40}{MJ}\ (\qty{11}{kWh}),\qquad \frac{39.8}{45}\approx\qty{0.88}{m^3}\]
\phtashikame
単位：kJ÷（kJ/kWh）＝kWh、MJ÷（MJ/m$^3$）＝m$^3$。\\
大きさ：ヒートポンプは、残りの $10.5-3.0=\qty{7.5}{kWh}$ を外の空気からくみ上げています。電気を効率0.60の発電所でつくるなら、燃料の熱は $3.0\div0.60=\qty{5.0}{kWh}$ で、ガス（\qty{11}{kWh}）の半分以下です。\\
極端な場合：$\mathrm{COP}_{\text{暖}}=1$ なら、電気のヒーターと同じ\qty{10.5}{kWh}になります。
\kotae （1）\qty{10.5}{kWh}（約\qty{38}{MJ}）（2）約\qty{3.0}{kWh}（3）約\qty{40}{MJ}、約\qty{0.88}{m^3}
\end{phreidai}

\begin{phreidai}[屋根の太陽光が1年につくる電気]\label{ph:熱-16:ex3}
定格の出力\qty{1.0}{kW}の太陽光の設備を家の屋根にのせます。設備利用率を0.15として、1年につくる電気を求めましょう。また、この設備の太陽電池の面積を、変換効率0.20、日光の強さ\qty{1000}{W/m^2}（よく晴れた昼）として見積もりましょう。
\phbamen
夜・雨・朝夕・冬の日をならした、1年の平均の発電量を知りたい。
\phmushi
設備利用率は、1年を通して一定の割合とみなします。パネルの温度による効率の変化や、パネルどうしのすき間は無視します。
\phzu
定格の出力 $P=\qty{1.0}{kW}$、設備利用率 0.15、1年 \qty{8760}{h}、変換効率 0.20、日光の強さ \qty{1000}{W/m^2}、面積 $A$。
\phriyu
設備利用率の式（⑤）を、1年につくる電気について解きます。面積は「日光の強さ×面積×効率＝定格の出力」から出します。
\phtateru
\[\text{1年の電気}=P\times\qty{8760}{h}\times0.15,\qquad A=\frac{P}{1000\times0.20}\]
\phtoku
\[1.0\times8760\times0.15\approx\qty{1314}{kWh},\qquad A=\frac{1000}{1000\times0.20}=\qty{5.0}{m^2}\]
\phtashikame
単位：kW×h＝kWh、W÷（W/m$^2$）＝m$^2$。\\
大きさ：1か月あたり約\qty{110}{kWh}で、熱-0 のやってみようの家（1か月\qty{260}{kWh}）の電気の4割ほどです。\\
極端な場合：設備利用率が1（いつも定格で発電）なら\qty{8760}{kWh}ですが、夜があるので、太陽光では決してそうなりません。
\kotae 約\qty{1300}{kWh}（年）、約\qty{5}{m^2}
\end{phreidai}

\begin{phreidai}[燃料電池と熱機関の上限をくらべる]\label{ph:熱-16:ex4}
水素と酸素から液体の水をつくる燃料電池の理論の効率（$\Delta H=\qty{-285.83}{kJ/mol}$、$\Delta G=\qty{-237.13}{kJ/mol}$）と、水素を燃やした熱を1000\,℃で受けとり、30\,℃で捨てる熱機関のカルノーの効率をくらべましょう。
\phbamen
同じ水素から電気をつくるのに、「燃やして熱機関で」と「燃料電池で」のどちらが、上限が高いかを知りたい。
\phmushi
どちらも、いちばんよい場合（可逆）を考えます。熱機関は、1000\,℃の熱源から熱を受けとるとします。
\phzu
燃料電池の理論の効率 $\eta_{\mathrm{FC}}=\Delta G/\Delta H$、カルノーの効率 $\eta_{\mathrm C}=1-T_{\mathrm L}/T_{\mathrm H}$、$T_{\mathrm H}=\qty{1273}{K}$、$T_{\mathrm L}=\qty{303}{K}$。
\phriyu
燃料電池は熱を通らないので、上限は $-\Delta G$（熱-10 ⑨）。熱機関の上限はカルノーの効率（熱-5 ⑤）。温度は K で入れます。
\phtateru
\[\eta_{\mathrm{FC}}=\frac{\Delta G}{\Delta H},\qquad \eta_{\mathrm C}=1-\frac{T_{\mathrm L}}{T_{\mathrm H}}\]
\phtoku
\[\eta_{\mathrm{FC}}=\frac{237.13}{285.83}\approx0.83,\qquad \eta_{\mathrm C}=1-\frac{303}{1273}\approx0.76\]
\phtashikame
大きさ：燃料電池のほうが、上限が高くなります。熱機関の上限が0.83になるのは、$1-303/T_{\mathrm H}=0.83$ を解いて、熱を受けとる温度が $T_{\mathrm H}\approx\qty{1780}{K}$（約1500\,℃）のときです。\\
本物では：燃料電池の電池1個の電圧は、理論の\qty{1.23}{V}より低い0.6〜0.8 V くらいで使うので、実際の効率は5割前後になります（熱-10）。上限が高いことと、実際に高いことは別です。
\kotae 燃料電池 約0.83、熱機関 約0.76
\end{phreidai}

\begin{phreidai}[効率の数を、同じ基準にそろえる]\label{ph:熱-16:ex5}
⑧の専門書の発電設備（発電端熱効率0.63、低発熱量で数えた値、所内率2\,\%）の、送電端の効率を、低発熱量と高発熱量の両方で求めましょう。燃料はメタン（高発熱量\qty{890.6}{kJ/mol}、低発熱量\qty{802.6}{kJ/mol}）とします。
\phbamen
同じ発電所の効率でも、数え方で数がちがいます。高発熱量・送電端で示されたほかの発電所の効率（熱-8 ⑧の約42\,\%）とくらべられるように、基準をそろえます。
\phmushi
燃料はすべてメタンとします。所内率は、運転のしかたによらず一定とします。
\phzu
電気（発電端）$P$、送電端の電気 $P_{\text{送}}=(1-0.02)P$、燃料の熱（低発熱量）$\dot Q_{\mathrm L}$、燃料の熱（高発熱量）$\dot Q_{\mathrm H}=\dot Q_{\mathrm L}\times890.6/802.6$。
\phriyu
効率は「電気÷燃料の熱」（熱-0 ②）。分子を送電端の電気に、分母を高発熱量で数えた熱に取りかえます。
\phtateru
\[\eta_{\text{送,低}}=\frac{P_{\text{送}}}{\dot Q_{\mathrm L}}=0.98\times0.63,\qquad \eta_{\text{送,高}}=\frac{P_{\text{送}}}{\dot Q_{\mathrm H}}=\eta_{\text{送,低}}\times\frac{802.6}{890.6}\]
\phtoku
\[\eta_{\text{送,低}}=0.98\times0.63\approx0.617,\qquad \eta_{\text{送,高}}=0.617\times0.901\approx0.556\]
\phtashikame
大きさ：基準をそろえると、0.63が0.56まで下がります。それでも、熱-8 ⑧の石炭火力（高発熱量・送電端で約0.42）より、ずっと高い効率です。\\
極端な場合：所内率が0、燃料が水素をふくまない（炭素だけ）なら、高発熱量と低発熱量は同じで、効率の数は変わりません。
\kotae 送電端の効率は、低発熱量で約0.62、高発熱量で約0.56
\end{phreidai}

%% ================================================================
\bksection{tsumazuki}{まちがえやすいところをチェック}
\begin{tsumazuki}{効率がいちばん高い方法}
\bkng{効率がいちばん高い発電の方法が、いちばんよい方法だ。}
\bkok{効率は物差しの1つ。燃料・CO$_2$・費用・出したいときに出せるか、なども見てくらべる。}
効率の分母（入れたもの）も方法でちがいます。太陽光の効率は約2割ですが、燃料は使いません（③）。
\end{tsumazuki}
\begin{tsumazuki}{ヒートポンプでエネルギーが増える？}
\bkng{成績係数3.5のヒートポンプは、電気のエネルギーを3.5倍に増やしている。}
\bkok{電気\qty{1}{kWh}と、外の空気からくみ上げた熱\qty{2.5}{kWh}をたした\qty{3.5}{kWh}を、お湯に入れている。}
エネルギーは保存しています（熱-5 ④）。得をしているのは、温度の低い熱をくみ上げて使っているからです。
\end{tsumazuki}
\begin{tsumazuki}{1 kWh は 1 kWh？}
\bkng{電気の\qty{1}{kWh}と、45\,℃のお湯の熱の\qty{1}{kWh}は、同じ価値だ。}
\bkok{エネルギーの量は同じでも、仕事にできる分（エクセルギー）がちがう。}
45\,℃のお湯の熱\qty{1}{kWh}のエクセルギーは、約\qty{0.05}{kWh}しかありません（⑥）。電気は、全部がエクセルギーです。
\end{tsumazuki}

%% ================================================================
\Needspace{12\baselineskip}\bksection{phgear}{工学ではこう使う}
\begin{phkougaku}[捨てる熱を使い回す：コージェネレーション]
熱-0 で見たように、発電所は燃料の熱の半分くらいを捨てています。その熱を、近くの工場の蒸気や、建物の暖房・給湯に使えば、燃料の熱をもっとむだなく使えます。このように、1つの燃料から電気と熱の両方をとり出すしくみを\textbf{コージェネレーション}（熱電併給）といいます。\bkidx{こーじぇねれーしょん}{コージェネレーション}工場・病院・ホテル・家庭用の燃料電池などで使われています。

ただし、熱は遠くへ運ぶと冷めてしまう（熱-11）ので、熱を使う場所の近くで発電する必要があります。また、③で見たように、熱には温度による「質」のちがいがあります。高い温度の熱でまず電気をつくり、残った低い温度の熱をお湯や暖房に使う、というように、温度の高い順に使っていくのが、エクセルギーをむだにしない使い方です。

技術者は、新しい設備を計画するとき、熱の流れの帯（熱-0）でエネルギーの行き先を数え、エクセルギーの帯（熱-6）で、どこで「質」が失われているかを数えます。損失の大きい所から改良していくのです。
\begin{phketa}[エネルギーを選ぶときの数]
\phketaline{電気\qty{1}{kWh}}{\qty{3.6}{MJ}}
\phketaline{1 kWh に要る燃料の熱（効率0.60／0.40）}{\qty{6.0}{MJ}／\qty{9.0}{MJ}}
\phketaline{燃やす所の CO$_2$（コンバインド／石炭）}{約\qty{330}{g}／約\qty{1000}{g}（1 kWh あたり）}
\phketaline{一生の CO$_2$（IPCC の中央値）}{石炭820、ガス490、風力・原子力 10 g ほど}
\phketaline{設備利用率（2024年の試算の前提）}{太陽光15〜18\,\%、風力約30\,\%}
\phketaline{燃料電池の理論の効率}{0.83（\qty{1.23}{V}）}
\phboxhead{答えが桁で合っているか}
1 kWh の電気をつくる燃料の熱は、熱機関なら\qty{3.6}{MJ}の約1.5〜3倍です。1 kWh あたりの CO$_2$ は、火力で数百 g、燃やさない方法で数十 g 以下です。kg と g、kWh と MJ を取りちがえると、1000倍や3.6倍ずれるので、単位を確かめます。
\end{phketa}
\end{phkougaku}
\phstation{エネルギー工学,電力系統,動力プラント,環境工学}{一次エネルギー、発電効率、高位発熱量・低位発熱量、設備利用率、ライフサイクルアセスメント、発電コスト、コージェネレーション、エクセルギー}

\begin{column}{エネルギーの単位いろいろ}
この本では、エネルギーを J で数えてきました。でも、エネルギーの統計や外国の資料には、ほかの単位も出てきます。電気は kWh（\qty{1}{kWh}＝\qty{3.6}{MJ}）、国際的なエネルギーの統計では「石油換算トン」（toe。石油\qty{1}{t}の熱にあたる量として、$10^7$\,kcal＝\qty{41.868}{GJ}と決めてある）、アメリカの資料では Btu（英熱量。水1ポンド（約\qty{454}{g}）を華氏1度（$\tfrac59$\,K）上げる熱で、約\qty{1055}{J}）がよく使われます。toe の kcal は、1\,kcal＝\qty{4.1868}{kJ} で、熱-1 の熱化学カロリー（\qty{4.184}{kJ}）とは少しちがう決め方です。単位がちがっても、2-10 の「単位つきの1をかける」換算で、すべて J にそろえられます。
\end{column}

\begin{furikaeri}
\begin{checklist}
\item \qty{1}{kWh}の電気に要る燃料の熱と、燃やして出る二酸化炭素の量を、効率と原子の数のバランスから見積もれる。\bkfurref{ph:熱-16:sec2}
\item 発電や給湯の方法を、効率・温度・設備利用率・一生で出る二酸化炭素など、いくつもの物差しでくらべられる。\bkfurref{ph:熱-16:sec9}
\item ヒートポンプや燃料電池が得をするわけを、エクセルギー（第二法則）で説明できる。\bkfurref{ph:熱-16:sec6}
\end{checklist}
\end{furikaeri}
\bkfukidashi{delta}{物理は「どれを選ぶか」までは決めてくれない。でも、くらべるための数を、正しく出すことはできるよ。この本の道具で、ニュースの数を確かめてみよう。}
